% !TeX root = ../main.tex

\chapter{Operational Safety Policy e Agent Guardrails}
\label{cap:safety_guardrails}

L'introduzione di layer agentici come MCP espone il sistema a rischi intrinseci legati ad allucinazioni, latenze e comportamenti non deterministici. Per mitigare tali rischi sono state introdotte specifiche \textit{Safety Policy} e vincoli architetturali, pensati per garantire che ogni azione degli agenti rimanga verificabile, confinata e sicura.

\section{Principio del Minimo Privilegio (Least Privilege)}
\label{sec:least_privilege}

In linea con le specifiche di progetto, l'architettura applica il principio del \textit{Least Privilege}: l'Intelligenza Artificiale è utilizzata come \emph{assistente operativo} e non può in alcun caso avere il controllo illimitato dell'infrastruttura.

Per implementare questo livello di sicurezza, i cicli di ragionamento dei due agenti sono confinati all'interno di un ambiente di esecuzione isolato gestito dall'orchestratore. Gli agenti non possiedono credenziali di root, né token diretti al cluster Kubernetes o al database InfluxDB. L'unica interfaccia verso il mondo reale è il server MCP, che agisce come \textit{chokepoint} di sicurezza applicando due regole fondamentali:

\begin{itemize}
    \item \textbf{Accesso Read-Only alla telemetria:} la telemetria operativa viene esposta esclusivamente in modalità di sola lettura. Gli strumenti di osservazione implementano sotto il cofano solo \texttt{query\_api} pre-compilate in Python. L'agente non ha facoltà di comporre query \textit{Flux} o \textit{SQL} arbitrarie, annullando matematicamente il rischio di \textit{injection}.
    \item \textbf{Divieto assoluto di azioni distruttive:} non è stato esposto alcun tool che permetta la compromissione dei dati primari. Azioni ad alto rischio come il \texttt{DROP} di un bucket, il \texttt{DELETE} di un ordine storico su InfluxDB o la terminazione dei processi core sono fisicamente impossibili per l'agente. Il codice del layer MCP (\texttt{drone\_mcp\_layer.py}) omette intenzionalmente qualsiasi implementazione logica di questi comandi.
\end{itemize}

Si crea così un perimetro di sicurezza per cui, anche nel caso di un'allucinazione logica in cui l'LLM tentasse autonomamente di cancellare dati operativi o distruggere il cluster, le richieste verrebbero rigettate per inesistenza del metodo, preservando l'integrità del sistema.

\section{Protezione contro Loop e Consumo Token}
\label{sec:protezione_loop}

Una vulnerabilità comune degli agenti è la possibilità che entrino in un loop di ragionamento infinito. Questo accade quando il modello subisce un'allucinazione e continua a invocare le stesse chiamate API senza giungere a una conclusione. Oltre a degradare il sistema, ciò porterebbe rapidamente all'esaurimento del budget di token.

L'\texttt{HealthAgent} dispone perciò di un contatore che traccia le chiamate effettuate nello stesso turno, impostato a 3 iterazioni: se l'agente non riesce a completare il ragionamento entro tale limite, il processo viene terminato forzatamente e viene restituito un messaggio predefinito. È inoltre presente un secondo controllo sul formato della risposta: se l'agente produce per due volte consecutive un output che non rispetta lo schema atteso, viene bloccato allo stesso modo. Questi due meccanismi limitano l'ingresso in cicli infiniti, mantenendo la spesa entro il budget stabilito.

\section{Protezione contro lo Scollegamento del Backend}
\label{sec:kill_switch_backend}

È stato implementato un interruttore di sicurezza per impedire lo spam di chiamate API verso l'LLM quando l'infrastruttura di backend va offline. Se il server MCP si disconnette, la funzione di recupero dello stato globale fallisce sistematicamente: senza protezioni, il loop continuerebbe a girare ogni pochi secondi inviando agli agenti contesti vuoti o dati obsoleti, costringendoli a prendere decisioni basate sul nulla e generando centinaia di invocazioni API a vuoto.

Per evitare questo spreco di token e risorse, ogni chiamata fallita verso il server incrementa un contatore dedicato. Al raggiungimento di una soglia prefissata scatta automaticamente un \textit{kill-switch} che congela l'attività degli agenti. Una volta attivata la sospensione globale, il sistema effettua un controllo periodico dello stato del backend; non appena il server MCP torna online, il contatore viene azzerato, l'interruttore di sicurezza si riarma e il flusso operativo riprende autonomamente la sua normale esecuzione.

\section{Idempotenza delle Azioni}
\label{sec:idempotenza_azioni}

In architetture di questo tipo, l'idempotenza è una proprietà di sicurezza fondamentale per garantire la robustezza del sistema e prevenire fallimenti catastrofici causati da allucinazioni o comportamenti ripetitivi degli LLM. Un'operazione è idempotente quando il suo effetto sullo stato del sistema è lo stesso, sia che venga eseguita una sola volta sia che venga eseguita più volte.

L'esempio più significativo di questa protezione è il meccanismo di scaling della flotta. Quando l'\texttt{HealthAgent} determina la necessità di espandere il numero di droni e invoca \texttt{scale\_drone\_deployment}, il codice in \texttt{drone\_mcp\_layer.py} esegue un'operazione di \textit{patching} sullo stato del cluster tramite la funzione \texttt{patch\_namespaced\_deployment\_scale}. Questa scelta garantisce l'idempotenza grazie a due fattori:

\begin{itemize}
    \item \textbf{Logica di stato desiderato (dichiarativa):} il tool accetta come parametro il numero di repliche desiderate (ad esempio 6). Se l'LLM entra in un loop e ripete cinque volte lo stesso comando, Kubernetes riceverà cinque richieste consecutive che dichiarano lo stesso stato target (6 Pod); il control-plane verificherà che il Deployment ha già raggiunto la quota di 6 repliche e ignorerà completamente le successive quattro chiamate, senza eseguire alcuna allocazione extra.
    \item \textbf{Prevenzione del collasso per esaurimento risorse:} se il tool MCP fosse stato progettato in modo \emph{incrementale} (aggiungendo un numero di droni a ogni chiamata), cinque chiamate consecutive dettate da un'allucinazione avrebbero inserito decine di droni fantasma nel cluster, saturando istantaneamente le risorse computazionali dei nodi, congestionando la rete MQTT e violando il tetto massimo di spesa.
\end{itemize}

\section{Human-In-The-Loop (HITL) ed Escalation}
\label{sec:hitl_escalation}

Il paradigma \textit{Human-In-The-Loop} (HITL) rappresenta l'ultimo livello di difesa contro allucinazioni o decisioni economicamente dannose generate autonomamente dall'Intelligenza Artificiale. Nel sistema, l'automazione pura viene interrotta non appena un'azione degli agenti supera i confini di sicurezza prestabiliti dalle policy, forzando un'escalation verso un supervisore umano.

\subsection*{La regola architetturale dei 6 droni}

All'interno del modulo \texttt{drone\_mcp\_layer.py}, la costante \texttt{POLICY\_LIMITS} fissa a \textbf{6} il numero massimo di droni allocabili in totale autonomia (\texttt{max\_drones\_auto}). Il flusso operativo è regolato dalla seguente policy:

\begin{itemize}
    \item \textbf{Sotto la soglia ($\le 6$):} l'\texttt{HealthAgent} ha pieno controllo dichiarativo e può invocare direttamente \texttt{scale\_drone\_deployment} per adattare il cluster Kubernetes ai picchi di carico.
    \item \textbf{Sopra la soglia ($> 6$):} qualora un workload maggiore richieda un dimensionamento superiore a 6 Pod, il layer MCP blocca l'esecuzione deterministica del comando considerandolo ad alto impatto economico. L'agente è obbligato a interrompere il proprio loop e a invocare lo strumento \texttt{request\_human\_approval}.
\end{itemize}

\subsection*{Il portale di approvazione}

L'intercettazione e la gestione delle escalation sono demandate a un servizio Flask indipendente (\texttt{human\_approval\_manager.py}) in ascolto sulla porta \texttt{5002}, che espone una dashboard web dedicata agli operatori umani.

Quando l'agente richiede un'escalation, il server MCP persiste la richiesta nel file \texttt{pending\_approvals.jsonl}. La dashboard legge questo registro e renderizza una scheda informativa per l'operatore. Ciascuna scheda mostra in modo trasparente l'identificativo univoco (\texttt{request\_id}), l'azione richiesta, il payload JSON (es. \texttt{\{"replicas": 8\}}) e, soprattutto, la giustificazione logica in linguaggio naturale fornita dall'LLM (\texttt{reason}).

L'operatore dispone di un diritto di veto assoluto tramite due interfacce: \textbf{Approve} convalida la transizione ed esegue istantaneamente il patch su Kubernetes impostando il parametro \texttt{force=True}, scavalcando i blocchi preventivi del codice; \textbf{Deny} contrassegna la richiesta come \texttt{rejected} e, al ciclo successivo, l'\texttt{HealthAgent} manterrà lo stato della flotta inalterato.
